\chapter{Validación y resultados}
\label{chap:validacion-resultados}

\section{Estrategia de validación}
\label{sec:estrategia-validacion}

La validación se ha organizado en tres niveles: pruebas de backend, validación de infraestructura y pruebas de operación sobre el entorno desplegado. Esta estrategia evita depender de una sola evidencia y permite demostrar tanto comportamiento funcional como capacidad real de despliegue.

En local se ejecutan pruebas automatizadas y flujos Postman/Newman contra una API viva. En infraestructura se usan \texttt{terraform fmt}, \texttt{terraform validate}, \texttt{terraform plan}, \texttt{terraform apply} y \texttt{terraform destroy}. En AWS se comprueba el estado de ECS, la migración puntual, la salud del endpoint, el bootstrap de administrador, el login y la existencia de recursos de observabilidad.

\section{Pruebas funcionales}
\label{sec:pruebas-funcionales}

La suite principal de pruebas automatizadas contiene 87 pruebas superadas tras Sprint 6. Estas pruebas cubren endpoints de salud, transportes, vehículos, conductores, autenticación, usuarios, reglas de estado y el contrato de correlación de peticiones. La validación automatizada se complementa con un flujo de smoke local basado en Postman/Newman, que arranca Docker Compose, prepara datos deterministas, ejecuta peticiones y limpia los datos generados.

Además, se verificaron manualmente el bootstrap del primer administrador y el login en el entorno cloud.

\section{Validación de infraestructura}
\label{sec:validacion-infraestructura}

La infraestructura se validó con Terraform antes y después del despliegue. Las acciones principales fueron:

\begin{itemize}
  \item Inicialización de backend remoto S3/DynamoDB.
  \item \texttt{terraform validate} correcto para el entorno \texttt{dev}.
  \item \texttt{terraform apply} real de la infraestructura \texttt{dev}.
  \item Publicación de imagen en ECR.
  \item Aplicación de migraciones EF Core mediante ECS task.
  \item Escalado del servicio ECS a una tarea.
  \item Activación del camino DNS/HTTPS con Route 53 y ACM.
  \item Creación de dashboard, alarmas y filtro métrico de errores.
  \item \texttt{terraform plan} con \path{-detailed-exitcode} sin cambios pendientes antes del cierre de coste.
  \item \texttt{terraform destroy} para eliminar recursos costosos del entorno \texttt{dev}.
\end{itemize}

\section{Resultados obtenidos}
\label{sec:resultados-obtenidos}

Los resultados técnicos obtenidos en la validación cloud son:

\begin{itemize}
  \item Endpoint de validación Sprint 6: DNS público del ALB registrado como evidencia operativa.
  \item Endpoint objetivo DNS/HTTPS: \url{https://api.dev.transitops.net}, gestionado por Route 53 y ACM cuando el entorno se recrea con \texttt{enable\_https = true}.
  \item Health check Sprint 6: \texttt{GET} \path{/api/v1/health/ready} devuelve \texttt{200}.
  \item Certificado ACM: validación DNS corregida en la cuenta \texttt{661000947340}; el certificado para \path{api.dev.transitops.net} llegó a estado \texttt{ISSUED} antes del destroy.
  \item Servicio ECS: \texttt{ACTIVE}, \texttt{desired = 1}, \texttt{running = 1}.
  \item Imagen ECR desplegada en Sprint 6: \path{sprint6-local-6bb910c}.
  \item Task definition desplegada: \path{transitops-dev-api:3}.
  \item Tarea ECS de migración: \path{345ae53d1340475aa94a826328404998}, finalizada con código \texttt{0}.
  \item Bootstrap de administrador: respuesta \texttt{201}.
  \item Login administrador: respuesta \texttt{200} con \acrshort{jwt}.
  \item Header de correlación: \texttt{X-Correlation-ID} presente en health, errores y login.
  \item CloudWatch Logs: entradas JSON filtrables por \texttt{State.CorrelationId}, incluyendo \path{sprint6-login-final}.
  \item Observabilidad: dashboard CloudWatch, log group de la API y 9 alarmas creadas por Terraform.
  \item SNS: topic no creado porque \texttt{alarm\_email} no se configuró.
  \item Cierre de coste: \texttt{terraform destroy} eliminó los recursos del entorno \texttt{dev}; el state quedó vacío.
\end{itemize}

Estos resultados demuestran que el sistema no solo compila, sino que se ejecuta sobre infraestructura cloud real y verificable, y que también puede destruirse de forma controlada para evitar gasto residual.

\section{Evidencia de cierre de coste}
\label{sec:evidencia-cierre-coste}

Tras la validación de Sprint 6 se ejecutó \texttt{terraform destroy}. El primer intento eliminó la mayoría de recursos y falló al borrar ECR porque el repositorio contenía imágenes. Se vació el repositorio con \texttt{aws ecr batch-delete-image} y se repitió el destroy, que finalizó correctamente. Después se verificó que:

\begin{itemize}
  \item \texttt{terraform state list} no devolvía recursos.
  \item No quedaban ALB, RDS, ECS ni ECR asociados a \path{transitops-dev}.
  \item El NAT Gateway aparecía en estado \texttt{deleted}.
  \item No quedaban certificados ACM ni registros Route 53 de \path{api.dev.transitops.net}.
  \item No quedaban log groups, dashboard, alarmas ni secretos del entorno \texttt{dev}.
\end{itemize}

Permanecen fuera de este destroy el dominio registrado \path{transitops.net}, su hosted zone pública y el backend remoto de Terraform, ya que son recursos de base necesarios para recreaciones posteriores y no forman parte del stack efímero de aplicación.

\section{Análisis de limitaciones}
\label{sec:analisis-limitaciones}

Las principales limitaciones son:

\begin{itemize}
  \item No existe interfaz gráfica; el sistema se consume mediante API, Postman o clientes equivalentes.
  \item El entorno desplegado es \texttt{dev}, no una arquitectura productiva multi-entorno completa.
  \item No se han implementado políticas de autoscaling ni pruebas de carga.
  \item La observabilidad cubre logs JSON, correlación, dashboard y alarmas básicas, pero no incluye trazas distribuidas ni APM completo.
  \item El pipeline de GitHub Actions está preparado y ampliado, pero todavía debe archivarse una ejecución completa desde GitHub como evidencia final de CI/CD.
  \item Algunas políticas IAM son prácticas para el sprint actual y pueden afinarse más en fases posteriores.
\end{itemize}
